\documentclass[11pt]{article}
\usepackage{amsmath,amssymb}
\begin{document}
\title{A direct nested construction of second-order AC subspaces}
\author{}\date{}\maketitle
Let $j=N/2$ and $c=j(j+1)/3$. Consider reflected first-order AC module
vectors
\[
 v_x=\sqrt{\frac{x+1}{2x+1}}\,|x\rangle+
     \sqrt{\frac{x}{2x+1}}\,|-x-1\rangle,
\]
where $|m\rangle=|D_N^{(j-m)}\rangle$ is the normalized spin/Dicke basis.
Choose positive module labels $x$ such that all magnetic support positions
of the selected $v_x$ are separated by at least three. In the canonical
reflection-pair construction, it suffices to choose its positive labels
spaced by three and omit any central block.

The vectors $v_x$ are orthonormal, have zero first magnetic moment, and
have second magnetic moment
\[
 h_x=x(x+1).
\]
Because every pair of occupied magnetic levels is separated by at least
three, all compressed ladder operators of shifts one and two vanish on
$\mathcal V=\operatorname{span}\{v_x\}$. Consequently a subspace
$\mathcal C\subseteq\mathcal V$ is second-order AC if and only if
\[
 P_{\mathcal C}\,\operatorname{diag}(h_x-c)\,P_{\mathcal C}=0.
\]
If $n_-,n_0,n_+$ count the negative, zero, and positive diagonal entries,
the largest dimension within this particular $\mathcal V$ is exactly
\[
 n_0+\min(n_-,n_+).
\]
Indeed, the inertia bound gives this upper bound. For existence, retain
all $v_x$ with $h_x=c$, pair disjoint labels with $h_x<c<h_y$, and use
\[
 u_{xy}=\sqrt{\frac{h_y-c}{h_y-h_x}}\,v_x+
         \sqrt{\frac{c-h_x}{h_y-h_x}}\,v_y.
\]
Disjoint supports make all cross moments diagonal, and every displayed
vector has the required second moment. Arbitrary relative phases in each
pair are permitted. This is an exact construction prescription, not a
classification of general first-order parents or an unrestricted
second-order optimality theorem.

\section*{An explicit eighteen-qubit example}
The canonical first-order parent for $N=18$ has dimension seven and includes
modules $x=2,5,8$. Their second moments are $6,30,72$, whereas $c=30$.
Thus
\[
 \mathcal C=\operatorname{span}\left\{v_5,
       \sqrt{\frac7{11}}v_2+\sqrt{\frac4{11}}v_8\right\}
\]
is second-order AC. In normalized Dicke coordinates, writing
$D_r=|D_{18}^{(r)}\rangle$, an orthonormal basis is
\[
 \psi_0=\sqrt{\frac6{11}}D_4+\sqrt{\frac5{11}}D_{15},
\]
\[
 \psi_1=\frac6{\sqrt{187}}D_1+
     \sqrt{\frac{21}{55}}D_7+
     \sqrt{\frac{14}{55}}D_{12}+
     \frac{4\sqrt2}{\sqrt{187}}D_{18}.
\]
The accompanying verifier checks orthonormality, all three compressed
first-spin matrices, all nine ordered second-spin products, and every
entry of all four two-qubit reduced cross matrices:
\[
 M_2(\psi_a)M_2(\psi_b)^\dagger=\frac{\delta_{ab}}3 I_3,
 \qquad a,b\in\{0,1\}.
\]
There are $4\times9=36$ scalar cross-reduction identities. The construction
is deliberately small: the unrestricted optimal dimension at $N=18$ is
four, so this example does not improve that catalogue entry.
\end{document}
